%=====================================================================
%                        AVISO 
%=====================================================================
% COPIE ESTE PROJETO PARA QUE POSSA EDITÁ-LO. ESSA OPÇÃO ESTÁ DISPONIVEL NO CANTO SUPERIOR ESQUERDO, EM 'MENU' --> 'COPY PROJECT'. SAIA DESTE LINK E COMECE A EDITAR O PROJETO-CÓPIA QUE ESTARÁ DISPONÍVEL NO SEU MENU!
%=====================================================================
%                        APROVEITE! 
%=====================================================================
\documentclass[final]{beamer}
% ====================
% Packages
% ====================
\usepackage[T1]{fontenc}
\usepackage[utf8]{luainputenc}
\usepackage[style=ieee]{biblatex}
\usepackage[sfdefault]{roboto}
\usepackage{placeins}
\usepackage[base]{babel}

\newif\ifazerosize
\ifazerosize
\usepackage[size=custom, width=84.1, height=118.9, scale=1.45]{beamerposter}
\else
\usepackage[size=custom, width=91.44, height=121.92, scale=1.45]{beamerposter}
\fi


\usetheme{gemini}
\usecolortheme{gt}
\usepackage{graphicx}

\usepackage{tikz}
\usepackage{pgfplots}
\pgfplotsset{compat=1.14}
\usepackage{anyfontsize}

\usepackage{fontawesome5}

\usepackage{lipsum}

\usepackage{tabularx}
\usepackage{booktabs}

\usepackage{siunitx}
\sisetup{detect-all}

\usepackage{amsfonts}
\usepackage{amsmath}
\usepackage{amssymb}

\hyphenpenalty=10000

\addbibresource{references.bib}

% ====================
% Lengths
% ====================
\newlength{\sepwidth}
\newlength{\colwidth}
\setlength{\sepwidth}{0.025\paperwidth}
\setlength{\colwidth}{0.4625\paperwidth}
\newcommand{\separatorcolumn}{\begin{column}{\sepwidth}\end{column}}
% ====================
% Title
% ====================
\title{TÍTULO TÍTULO TÍTULO TÍTULO TÍTULO TÍTULO}

\author{A. Autor, B. Autor}

\institute[shortinst]{Instituição Instituição Instituição}
% ====================
% Footer (optional)
% ====================
\footercontent{
    \href{mailto:email@email}{\faEnvelope\hspace{0.5ex} email@email}}
% ====================
% Logo
% ====================
\logoright{\includegraphics[height=15cm]{logos/logo-1.png}}
%coloque o logo da sua instituição aqui
\logoleft{\includegraphics[height=15cm]{logos/logo_inst.png}}
% ====================




% ====================
% Body
% ====================
\begin{document}
\begin{frame}[t]
\begin{columns}[t]
\separatorcolumn
\begin{column}{\colwidth}
% ====================
% Block 1
% ====================
\begin{block}{Introdução}
\lipsum[1] \cite{nussenzveig2013curso}
\end{block}
% ====================
% Block 2
% ====================
\begin{block}{Objetivos}
Texto texto texto texto texto.
\end{block}
% ====================
% Block 3
% ====================
\begin{block}{Materiais e Métodos (Metodologia)}
\lipsum[1]
\begin{equation}
    \label{dft_1}
    \frac{-\hbar}{2m} \frac{d^2}{dr^2} \psi(\boldsymbol{r}) + V(\boldsymbol{r})\psi (\boldsymbol{r}) = E \psi(\boldsymbol{r}),
\end{equation}
\end{block}
% ====================
\end{column}
\separatorcolumn
\begin{column}{\colwidth}
% ====================
% Block 4
% ====================
\begin{block}{Resultados e Discussões}
\begin{table}[htb]
    \centering
    \caption{Exemplo tabela.}
    \small
    \label{TABELA1}
    \begin{tabular}{l|c c c c c c }
        \hline
        A & 1 & 5 & 9 & 13 \\
        B & 2 & 6 & 10 & 14 \\
        C & 3 & 7 & 11 & 15 \\
        D & 4 & 8 & 12 & 16 \\ 
        \hline
    \end{tabular}
\end{table}
\FloatBarrier

\begin{figure}
\centering
\includegraphics[height=20cm]{figures/exemplo.png}
\caption{Exemplo de figura.}
\label{FIGURA}
\end{figure}
\end{block}
% ====================
% Block 5
% ====================
\begin{block}{Conclusões}
\lipsum[1]       
\end{block}
% ====================
% Block 6
% ====================
\begin{block}{Agradecimentos}

\begin{figure}[b]
\centering
\includegraphics[width =8cm]{logos/Logo_Unesp.svg.png}
\includegraphics[width=6cm]{logos/brasaoigce.png}
\vspace{0.5cm}
\end{figure}

\begin{figure}[b]
\centering
\includegraphics[width =6cm]{logos/Fapesp.png}
\includegraphics[width =6cm]{logos/capes.jpg}
\includegraphics[width=6cm]{logos/cnpq.jpg}
\vspace{0.5cm}
\end{figure}

\end{block}
% ====================
% Block 7
% ====================
\begin{block}{Referências}

%\nocite{Vaswani2017AttentionNeed}
\AtNextBibliography{\small}
\printbibliography

\end{block}
% ====================    
\end{column}
\separatorcolumn
\end{columns}
\end{frame}

\end{document}
% ====================